\appendixitem{REPRESENTATIONS OF COMPACT GROUPS}

Let \(G\) be a compact topological group, and let \(\rho : G \to \mathbf{GL}(W)\) be a continuous representation of \(G\) on a finite dimensional complex vector space. We shall say that \(\rho\) is irreducible of real, complex or quaternionic type if this is the case for the \(\mathbf{C}^{(G)}\) -module \(W\) (relative to the algebra \(A = \mathbf{R}^{(G)}\) ). Let \(H\) be a separating positive hermitian form on \(W\) , invariant under \(G\) .

PROPOSITION 3. Assume that \(\rho\) is irreducible.

\begin{enumerate}
\def\labelenumi{\alph{enumi})}
\item
  The representation \(\rho\) is of real type if and only if there exists a non-zero symmetric bilinear form B on W, invariant under G. In this case the form B is separating; the set V of \(w \in W\) such that \(\mathrm{H}(w, x) = \mathrm{B}(w, x)\) for all \(x \in W\) is an \(\mathbf{R}\)-structure on W invariant under G.
\item
  The representation \(\rho\) is of complex type if and only if there exists no non-zero bilinear form on W invariant under G.
\item
  The representation \(\rho\) is of quaternionic type if and only if there exists a non-zero alternating bilinear form on W, invariant under G; such a form is necessarily separating.
\end{enumerate}

For \(\theta \in \mathrm{Hom}_{\mathbf{C}^{(\mathrm{G})}}(\mathrm{W},\overline{\mathrm{W}})\) and \(x,y\in \mathrm{W}\) , put \(\mathrm{B}_{\theta}(x,y) = \mathrm{H}(\theta x,y)\) . Then \(\mathrm{B}_{\theta}\) is a bilinear form on \(\mathrm{W}\) , invariant under \(\mathrm{G}\) , and separating if \(\theta\) is non-zero. Denote by \(\mathcal{B}(\mathrm{W})^{\mathrm{G}}\) the space of bilinear forms on \(\mathrm{W}\) invariant under \(\mathrm{G}\) ; the map \(\theta \mapsto \mathrm{B}_{\theta}\) from \(\mathrm{Hom}_{\mathbf{C}^{(\mathrm{G})}}(\mathrm{W},\overline{\mathrm{W}})\) to \(\mathcal{B}(\mathrm{W})^{\mathrm{G}}\) is an isomorphism of \(\mathbf{C}\) -vector spaces. This implies, in particular, assertion \(b\) ).

Let \(\theta\) be a \(\mathbf{C}^{(\mathrm{G})}\) -isomorphism from W to \(\overline{\mathrm{W}}\) such that \(\theta \circ \theta = \alpha_{\mathrm{W}}\) , with \(\alpha \in \{-1, +1\}\) (Prop. 2); since \(\mathcal{B}(\mathrm{W})^{\mathrm{G}}\) is of dimension 1, there exists \(\varepsilon \in \mathbf{C}\) such that

\[
\mathrm{B} _ {\theta} (y, x) = \varepsilon \mathrm{B} _ {\theta} (x, y) \quad \text { for   all } x, y \text { in } \mathrm{W}.
\]

Iterating, we obtain \(\mathrm{B}_{\theta}(y,x)=\varepsilon\mathrm{B}_{\theta}(x,y)=\varepsilon^{2}\mathrm{B}_{\theta}(y,x)\) , so \(\varepsilon^{2}=1\) and \(\varepsilon\in\{-1,+1\}\) . Moreover, for x in W,

\[
\mathrm{H} (\theta x, \theta x) = \mathrm{B} _ {\theta} (x, \theta x) = \varepsilon \mathrm{B} _ {\theta} (\theta x, x) = \varepsilon \mathrm{H} (\theta \circ \theta (x), x) = \varepsilon \alpha \mathrm{H} (x, x)
\]

so \(\varepsilon\alpha > 0\) since H is positive, that is, \(\varepsilon = \alpha\) . Assertions a) and c) now follow from Prop. 2.

Denote by dg the Haar measure of total mass 1 on G.

Lemma 1. Let \(W^{G}\) be the subspace of W consisting of the elements invariant under G. The endomorphism \(\int_{G}\rho(g)dg\) of W is a projection with image \(W^{G}\) , compatible with the operations of G. In particular,

\[
\dim \mathrm{W} ^ {\mathrm{G}} = \int_ {\mathrm{G}} \operatorname{Tr} \rho (g) d g.
\]

Put \(p = \int_{\mathrm{G}} \rho(g) dg\) ; for \(h \in \mathrm{G}\) ,

\[
\rho (h) \circ p = \int_ {\mathrm{G}} \rho (h g) d g = \int_ {\mathrm{G}} \rho (g) d g = p
\]

and similarly \(p \circ \rho(h) = p\) . Thus, p is compatible with the operations of G, and its image is contained in \(W^{G}\) . If \(w \in W^{G}\) , we have \(p(w) = \int_{\mathbf{G}} \rho(g) w \, dg = w\) , hence the lemma.

Lemma 2. Let \(u\) be an endomorphism of a finite dimensional vector space \(E\) over a field \(K\) . Then

\[
\operatorname{Tr} u ^ {2} = \operatorname{Tr} \mathbf {S} ^ {2} (u) - \operatorname{Tr} \Lambda^ {2} (u).
\]

Let \(\chi_u(\mathrm{X}) = \prod_{i=1}^{n} (\mathrm{X} - \alpha_i)\) be a decomposition of the characteristic polynomial of \(u\) into linear factors in a suitable extension of K. We have \(\operatorname{Tr} u^2 = \sum_i \alpha_i^2\) , \(\operatorname{Tr} \Lambda^2(u) = \sum_{i < j} \alpha_i \alpha_j\) , \(\operatorname{Tr} \mathbf{S}^2(u) = \sum_{i \leq j} \alpha_i \alpha_j\) (cf.~Algebra, Chap. VII, §5, no. 5, Cor. 3), hence the result.

PROPOSITION 4. Assume that \(\rho\) is irreducible. Then, \(\rho\) is of real (resp. complex, resp. quaternionic) type if and only if the integral \(\int_{\mathrm{G}} \mathrm{Tr} \rho(g^2) dg\) is equal to 1 (resp. 0, resp. -1).

Denote by \(\check{\rho}\) the contragredient representation of \(\rho\) on \(W^{*}\) (defined by \(\check{\rho}(g)={}^{t}\rho(g^{-1})\) ). Applying Lemma 2 to \(\check{\rho}(g)\) and integrating over G gives

\[
\int_ {\mathrm{G}} \operatorname{Tr} \rho (g ^ {2}) d g = \int_ {\mathrm{G}} \operatorname{Tr} ^ {t} \rho (g ^ {- 2}) d g = \int_ {\mathrm{G}} \operatorname{Tr} \mathbf {S} ^ {2} (\check {\rho} (g)) d g - \int_ {\mathrm{G}} \operatorname{Tr} \Lambda^ {2} (\check {\rho} (g)) d g
\]

hence, by Lemma 1,

\[
\int_ {\mathrm{G}} \mathrm{Tr} \rho (g ^ {2}) d g = \dim (\mathbf {S} ^ {2} \mathrm{W} ^ {*}) ^ {\mathrm{G}} - \dim (\Lambda^ {2} \mathrm{W} ^ {*}) ^ {\mathrm{G}}.
\]

But \(S^{2}W^{*}\) (resp. \(\Lambda^{2}W^{*}\) ) can be identified with the space of symmetric (resp. alternating) bilinear forms on W. Thus, the proposition follows immediately from Prop. 3.

\specialchapter{EXERCISES}

\exercisesection{\S~1}

\begin{enumerate}
\def\labelenumi{\arabic{enumi})}
\tightlist
\item
  Let \(\mathrm{G}\) be a finite dimensional, connected, commutative complex Lie group, and let \(\mathrm{V}\) be its Lie algebra.
\end{enumerate}

\begin{enumerate}
\def\labelenumi{\alph{enumi})}
\item
  The map \(\exp_{\mathrm{G}}: \mathrm{V} \to \mathrm{G}\) is a surjective homomorphism; its kernel \(\Gamma\) is a discrete subgroup of \(\mathrm{V}\) .
\item
  G is compact if and only if \(\Gamma\) is a lattice in V; then G is said to be a complex torus.
\item
  Let \(\Gamma\) be the discrete subgroup of \(\mathbf{C}^2\) generated by the elements \(e_1 = (1,0), e_2 = (0,1), e_3 = (\sqrt{2},i)\) ; put \(\mathrm{G} = \mathrm{C}^2 / \Gamma\) , \(\mathrm{H} = (\Gamma + \mathrm{C}e_1) / \Gamma\) . Show that H is isomorphic to \(\mathbf{C}^*\) , that G/H is a complex torus of dimension 1, but that G contains no non-zero complex torus.
\item
  Every finite dimensional connected compact complex Lie group is a complex torus (cf.~Chap. III, §6, no. 3, Prop. 6).
\end{enumerate}

\begin{enumerate}
\def\labelenumi{\arabic{enumi})}
\setcounter{enumi}{1}
\tightlist
\item
  Let H be the set of complex numbers \(\tau\) such that \(\mathcal{I}(\tau) > 0\) .
\end{enumerate}

\begin{enumerate}
\def\labelenumi{\alph{enumi})}
\item
  Show that an analytic law of left operation of the discrete group \(\mathbf{SL}(2,\mathbf{Z})\) on H can be defined by putting \(\gamma\tau=\frac{a\tau+b}{c\tau+d}\) for all \(\gamma=\begin{pmatrix}a&b\\ c&d\end{pmatrix}\in\mathbf{SL}(2,\mathbf{Z})\) and \(\tau\in H\) .
\item
  For \(\tau\in H\) , denote by \(T_{\tau}\) the complex torus \(\mathbf{C}/(\mathbf{Z}+\mathbf{Z}\tau)\) . Show that the map \(\tau\mapsto T_{\tau}\) induces by passage to the quotient a bijective map from the set \(\mathrm{H}/\mathrm{SL}(2,\mathbf{Z})\) to the set of isomorphism classes of complex tori of dimension one.
\end{enumerate}

\begin{enumerate}
\def\labelenumi{\arabic{enumi})}
\setcounter{enumi}{2}
\item
  Let G be an integral subgroup of \(\mathbf{O}(n)\) such that the identity representation of G on \(\mathbf{R}^{n}\) is irreducible. Show that G is closed (write G in the form \(K \times N\) , where K is compact and N is commutative, and prove that \(\dim \overline{N} \leq 1\) ).
\item
  Show that the conditions in Prop. 3 (no. 3) are equivalent to each of the following conditions:
\end{enumerate}

(ii') The group \(\operatorname{Ad}(\mathbf{G})\) is relatively compact in \(\operatorname{Aut}(\mathbf{L}(\mathbf{G}))\) .

(ii'\,') The group \(\operatorname{Ad}(\mathbf{G})\) is relatively compact in \(\operatorname{End}(\mathbf{L}(\mathbf{G}))\) .

\begin{enumerate}
\def\labelenumi{(\alph{enumi})}
\setcounter{enumi}{21}
\tightlist
\item
  Every neighbourhood of the identity element e of G contains a neighbourhood of e stable under inner automorphisms (cf.~Integration, §3, no. 1, Prop. 1).
\end{enumerate}

\begin{enumerate}
\def\labelenumi{\arabic{enumi})}
\setcounter{enumi}{4}
\tightlist
\item
  Let A be the closed subgroup of \(\mathbf{GL}(3,\mathbf{R})\) consisting of the matrices \((a_{ij})\) such that \(a_{ij}=0\) for i\textgreater j, \(a_{ii}=1\) for \(1\leq i\leq3\) , \(a_{12}\in \mathbf{Z}\) , \(a_{23}\in \mathbf{Z}\) ; let \(\theta\in \mathbf{R}-\mathbf{Q}\) and let B be the subgroup of A consisting of the matrices \((a_{ij})\) such that \(a_{12}=a_{23}=0\) and \(a_{13}\in\theta\mathbf{Z}\) , and let G=A/B.
\end{enumerate}

\begin{enumerate}
\def\labelenumi{\alph{enumi})}
\item
  Show that G is a Lie group of dimension one, whose Lie algebra is compact.
\item
  Show that \(\mathrm{C}(\mathrm{G}) = \overline{\mathrm{D}(\mathrm{G})} = \mathrm{G}_{0}\) , and that \(G_{0}\) is the largest compact subgroup of G.
\item
  Show that G is not the semi-direct product of \(G_{0}\) with any subgroup.
\end{enumerate}

\begin{enumerate}
\def\labelenumi{\arabic{enumi})}
\setcounter{enumi}{5}
\item
  Show that a (real) Lie algebra is compact if and only if it admits a basis \((e_{\lambda})_{\lambda\in\mathrm{L}}\) for which the structure constants \(\gamma_{\lambda\mu\nu}\) are anti-symmetric in \(\lambda,\mu,\nu\) (that is, they satisfy \(\gamma_{\lambda\mu\nu}=-\gamma_{\mu\lambda\nu}=-\gamma_{\lambda\nu\mu}\) ).
\item
  An involutive Lie algebra is a (real) Lie algebra \(\mathfrak{g}\) equipped with an automorphism s such that \(s \circ s = 1_{\mathfrak{g}}\) . Denote by \(\mathfrak{g}^{+}\) (resp. \(\mathfrak{g}^{-}\) ) the eigenspace of \(\mathfrak{g}\) relative to the eigenvalue +1 (resp. -1) of s.
\end{enumerate}

\begin{enumerate}
\def\labelenumi{\alph{enumi})}
\item
  Show that \(\mathfrak{g}^{+}\) is a subalgebra of \(\mathfrak{g}\) and that \(\mathfrak{g}^{-}\) is a \(\mathfrak{g}^{+}\) -module; we have \([\mathfrak{g}^{-},\mathfrak{g}^{-}] \subset \mathfrak{g}^{+}\) , and \(\mathfrak{g}^{+}\) and \(\mathfrak{g}^{-}\) are orthogonal with respect to the Killing form.
\item
  Show that the following conditions are equivalent:
\end{enumerate}

\begin{enumerate}
\def\labelenumi{(\roman{enumi})}
\item
  The \(\mathfrak{g}^{+}\) -module \(\mathfrak{g}^{-}\) is simple.
\item
  The algebra \(\mathfrak{g}^{+}\) is maximal among the subalgebras of \(\mathfrak{g}\) distinct from \(\mathfrak{g}\). If these conditions are satisfied, and if \(\mathfrak{g}^{+}\) contains no non-zero ideal of \(\mathfrak{g}\), the involutive Lie algebra \((\mathfrak{g}, s)\) is said to be irreducible.
\end{enumerate}

\begin{enumerate}
\def\labelenumi{\alph{enumi})}
\setcounter{enumi}{2}
\tightlist
\item
  Assume that \(\mathfrak{g}\) is semi-simple. Show that the following conditions are equivalent:
\end{enumerate}

\begin{enumerate}
\def\labelenumi{(\roman{enumi})}
\item
  The only ideals of \(\mathfrak{g}\) stable under \(s\) are \(\{0\}\) and \(\mathfrak{g}\) ;
\item
  \(\mathfrak{g}\) is either simple or the sum of two simple ideals interchanged by \(s\) .
\end{enumerate}

Show that these conditions are satisfied when \((\mathfrak{g}, s)\) is irreducible (observe that \(\mathfrak{g}\) is the direct sum of ideals stable under s and satisfying (ii)).

\begin{enumerate}
\def\labelenumi{\alph{enumi})}
\setcounter{enumi}{3}
\item
  Assume from now on that \(\mathfrak{g}\) is compact semi-simple. Show that the involutive Lie algebra \((\mathfrak{g}, s)\) is irreducible if and only if s is different from the identity and \((\mathfrak{g}, s)\) satisfies the equivalent conditions in c) (let p be a \(\mathfrak{g}^{+}\) -submodule of \(\mathfrak{g}^{-}\) and q its orthogonal complement with respect to the Killing form; observe that \([p, q] = 0\) and deduce that \(p + [p, p]\) is an ideal of \(\mathfrak{g}\)).
\item
  Prove that \(\mathfrak{g}\) is the direct sum of a family \((\mathfrak{g}_{i})_{0\leq i\leq n}\) of ideals stable under s, such that \(\mathfrak{g}_{0}\) is fixed by s and such that the involutive algebras \((\mathfrak{g}_{i},s|\mathfrak{g}_{i})\) are irreducible for \(1\leq i\leq n\) .
\end{enumerate}

\begin{enumerate}
\def\labelenumi{\arabic{enumi})}
\setcounter{enumi}{7}
\tightlist
\item
  Let G be a compact semi-simple Lie group, u an automorphism of G of order two, K the identity component of the set of fixed points of u, and X the manifold G/K (the homogeneous space X is then said to be symmetric).
\end{enumerate}

\begin{enumerate}
\def\labelenumi{\alph{enumi})}
\item
  Show that, if G is almost simple, then K is maximal among the connected closed subgroups of G distinct from G; in other words (Chap. III, §3, Exerc. 8), the operation of G on X is primitive. The symmetric space X is then said to be irreducible.
\item
  Assume that the manifold X is simply-connected; show that it is then isomorphic to a product of manifolds each of which is isomorphic either to a Lie group or to an irreducible symmetric space.
\end{enumerate}

\begin{enumerate}
\def\labelenumi{\arabic{enumi})}
\setcounter{enumi}{8}
\tightlist
\item
  Let \(\mathfrak{a}\) be a real or complex Lie algebra, and let \(G\) be a compact subgroup of \(\operatorname{Aut}(\mathfrak{a})\) . Prove that \(\mathfrak{a}\) has a Levi subalgebra (Chap. I, §6, no. 8, Def. 7) stable under \(G\) (reduce to the case in which the radical of \(\mathfrak{a}\) is abelian, and use Integration, Chap. VII, §3, no. 2, Lemma 2).
\end{enumerate}

\exercisesection{\S~2}

\begin{enumerate}
\def\labelenumi{\arabic{enumi})}
\tightlist
\item
  Let \(\mathbf{G}\) be a connected compact Lie group, \(g\) an element of \(\mathbf{G}\) .
\end{enumerate}

\begin{enumerate}
\def\labelenumi{\alph{enumi})}
\item
  Show that there exists an integer \(n \geq 1\) such that the centralizer \(Z(g^{n})\) is connected (prove that, for suitable n, the closed subgroup generated by \(g^{n}\) is a torus).
\item
  Assume that the dimension of \(\operatorname{Ker}(\operatorname{Ad}g^n - 1)\) is independent of \(n\) ( \(n \geq 1\) ). Prove that \(Z(g)\) is connected.
\item
  If \(g^n\) is regular for all \(n \geq 1\) , \(Z(g)\) is connected.
\end{enumerate}

\begin{enumerate}
\def\labelenumi{\arabic{enumi})}
\setcounter{enumi}{1}
\item
  Show that every connected compact Lie group G is the semi-direct product of its derived group by a torus (observe that, if T is a maximal torus of G, then \(D(G) \cap T\) is a torus).
\item
  Let G be a compact Lie group, \(\mathfrak{g}\) its Lie algebra, \(\mathfrak{s}\) a vector subspace of \(\mathfrak{g}\) such that, for \(x, y, z\) in \(\mathfrak{s}\), we have \([x, [y, z]] \in \mathfrak{s}\) . Let \(\mathcal{L}\) be the set of commutative subalgebras of \(\mathfrak{g}\) contained in \(\mathfrak{s}\). Show that the identity component of the stabilizer of \(\mathfrak{s}\) in G operates transitively on the set of maximal elements of \(\mathcal{L}\) (argue as in the proof of Th. 1).
\item
  Let G be a connected compact Lie group, T a maximal torus of G and S a sub-torus of T. Denote by \(\Sigma\) (resp. F) the stabilizer (resp. the fixer) of S in \(W_{G}(T)\) .
\end{enumerate}

\begin{enumerate}
\def\labelenumi{\alph{enumi})}
\item
  Prove that the group \(\mathrm{N_G(S) / Z_G(S)}\) is isomorphic to the quotient \(\Sigma /\mathrm{F}\) .
\item
  Let H be a connected closed subgroup of G, such that S is a maximal torus of H. Show that every element of \(W_{\mathrm{H}}(\mathrm{S})\) , considered as an automorphism of S, is the restriction to S of an element of \(\Sigma\) .
\end{enumerate}

\begin{enumerate}
\def\labelenumi{\arabic{enumi})}
\setcounter{enumi}{4}
\tightlist
\item
  Let \(G\) be a connected compact Lie group and \(S\) a torus of \(G\) . Show that the following conditions are equivalent:
\end{enumerate}

\begin{enumerate}
\def\labelenumi{(\roman{enumi})}
\item
  S is contained in a unique maximal torus;
\item
  \(Z_{G}(S)\) is a maximal torus;
\item
  S contains a regular element.
\end{enumerate}

\begin{enumerate}
\def\labelenumi{\arabic{enumi})}
\setcounter{enumi}{5}
\tightlist
\item
  Let G be a connected compact Lie group, T a maximal torus of G, \(\mathfrak{g}\) (resp. \(\mathfrak{t}\)) the Lie algebra of G (resp. T), and \(i : \mathfrak{t} \to \mathfrak{g}\) the canonical injection. Prove that the map \(^{t}i : \mathfrak{g}^{*} \to \mathfrak{t}^{*}\) induces by passage to the quotient a homeomorphism from \(\mathfrak{g}^{*}/G\) to \(\mathfrak{t}^{*}/W_{G}(T)\) .
\end{enumerate}

¶ 7) Let X and Y be two real manifolds of class \(C^{r}\) ( \(1 \leq r \leq \omega\) ), separated, connected and of dimension n; let \(f : X \to Y\) be a proper morphism of class \(C^{r}\) , equipped with an orientation F (Differentiable and Analytic Manifolds, Results, 10.2.5).

\begin{enumerate}
\def\labelenumi{\alph{enumi})}
\item
  Show that there exists a unique real number d such that \(\int_{X} f^{*}\alpha = d\int_{Y}\alpha\) for every twisted differential form \(\alpha\) of degree n on Y of class \(C^{1}\) and of compact support (use Differentiable and Analytic Manifolds, Results, 11.2.4).
\item
  Let \(y \in \mathbf{Y}\) be such that \(f\) is étale at every point of \(f^{-1}(y)\) . For \(x \in f^{-1}(y)\) , put \(\nu_x(f) = 1\) (resp. \(\nu_x(f) = -1\) ) if the maps \(\mathrm{F}_x\) and \(\tilde{f}_x\) (Differentiable and Analytic Manifolds, Results, 10.2.5, Example b)) from \(\mathrm{Or}(\mathrm{T}_x(\mathrm{X}))\) to \(\mathrm{Or}(\mathrm{T}_y(\mathrm{Y}))\) coincide (resp. are opposite). Prove that \(d = \sum_{x \in f^{-1}(y)} \nu_x(f)\) and, in particular, that \(d \in \mathbf{Z}\) .
\end{enumerate}

We call d the degree of f and denote it by \(\deg f\) . If X = Y and X is orientable, we can take for F the orientation that preserves the orientation of X.

\begin{enumerate}
\def\labelenumi{\alph{enumi})}
\setcounter{enumi}{2}
\item
  Show that, if \(\deg f \neq 0\) then \(f\) is surjective.
\item
  If there exists \(x \in X\) such that \(f^{-1}(f(x)) = \{x\}\) and such that f is étale at x, then f is surjective.
\end{enumerate}

\begin{enumerate}
\def\labelenumi{\arabic{enumi})}
\setcounter{enumi}{7}
\tightlist
\item
  Let G be a connected compact Lie group, dg the normalized Haar measure on G.
\end{enumerate}

\begin{enumerate}
\def\labelenumi{\alph{enumi})}
\tightlist
\item
  Let \(f: \mathrm{G} \to \mathrm{G}\) be a morphism of manifolds; for \(g \in \mathrm{G}\) , the differential \(\mathrm{T}_g(f)\) can be identified by left translations with a linear map from \(\mathrm{L}(\mathrm{G})\) to \(\mathrm{L}(\mathrm{G})\) . Prove the formulas
\end{enumerate}

\[
\deg f. \int_ {\mathrm{G}} \varphi (g) d g = \int_ {\mathrm{G}} \varphi (f (g)) \det \mathrm{T} _ {g} (f) d g
\]

for every integrable (complex-valued) function \(\varphi\) on G, and

\[
\deg f = \int_ {\mathrm{G}} \det \mathrm{T} _ {g} (f) d g.
\]

\begin{enumerate}
\def\labelenumi{\alph{enumi})}
\setcounter{enumi}{1}
\tightlist
\item
  Let \(\psi_k: \mathbf{G} \to \mathbf{G}\) be the map \(g \mapsto g^k\) . Prove the formula
\end{enumerate}

\[
\deg \psi_ {k} = \int_ {\mathrm{G}} \det (1 + \operatorname{Ad} g + \dots + (\operatorname{Ad} g) ^ {k - 1}) d g.
\]

\begin{enumerate}
\def\labelenumi{\alph{enumi})}
\setcounter{enumi}{2}
\tightlist
\item
  Let \(\mathrm{T}\) be a maximal torus of \(\mathrm{G}\) , and \(\mathrm{T}_r\) the set of regular elements of \(\mathrm{T}\) . Show that \(\exp_{\mathrm{G}}^{-1}(\mathrm{T}_r) \subset \mathrm{L}(\mathrm{T})\) and \(\psi_k^{-1}(\mathrm{T}_r) \subset \mathrm{T}_r\) for all \(k \geq 1\) .
\end{enumerate}

\(d\)) Show that \(\deg \psi_{k} = k^{\dim (\mathrm{T})}\) ; deduce the equality

\[
\int_ {\mathrm{G}} \det (1 + \operatorname{Ad} g + \dots + (\operatorname{Ad} g) ^ {k - 1}) d g = k ^ {\dim (\mathrm{T})}.
\]

\begin{enumerate}
\def\labelenumi{\arabic{enumi})}
\setcounter{enumi}{8}
\tightlist
\item
  This exercise is devoted to another proof of Th. 2. Let G be a connected compact Lie group, T a maximal torus of G, N its normalizer. Denote by \(G \times^{N} T\) the quotient manifold of \(G \times T\) by N for the operation defined by \((g, t).n = (gn, n^{-1}tn)\) (Differentiable and Analytic Manifolds, Results, 6.5.1).
\end{enumerate}

\begin{enumerate}
\def\labelenumi{\alph{enumi})}
\item
  Show that the morphism \((g, t) \mapsto gtg^{-1}\) from \(G \times T\) to G defines by passage to the quotient an analytic morphism \(f : G \times^{N} T \to G\) .
\item
  Let \(\theta\) be an element of T whose set of powers is dense in T (General Topology, Chap. VII, §1, no. 3, Cor. 2 of Prop. 7). Show that \(f^{-1}(\theta) = \{x\}\) , where \(x\) is the class of \((e,\theta)\) in \(\mathrm{G} \times^{\mathrm{N}}\mathrm{T}\) , and that \(f\) is étale at \(x\) .
\item
  Conclude from Exerc. 7 d) that \(f\) is surjective, and deduce another proof of Th. 2.
\end{enumerate}

\begin{enumerate}
\def\labelenumi{\arabic{enumi})}
\setcounter{enumi}{9}
\tightlist
\item
  * In this exercise, we use the following result from algebraic topology (Lefschetz's formula \(^{11}\) ): let X be a finite dimensional compact manifold, and f a morphism from X to itself. Assume that the set F of points x of X such that \(f(x) = x\) is finite, and that for all \(x \in F\) the number \(\delta(x) = \det(1 - \mathrm{T}_{x}(f))\) is non-zero. If \(\mathrm{H}^{i}(f)\) denotes the endomorphism of the vector space \(\mathrm{H}^{i}(\mathrm{X}, \mathbf{R})\) induced by f (for \(i \geq 0\) ), then
\end{enumerate}

\[
\sum_ {x \in \mathrm{F}} \delta (x) / | \delta (x) | = \sum_ {i \geq 0} (- 1) ^ {i} \operatorname{Tr} \mathrm{H} ^ {i} (f).
\]

Let G be a connected compact Lie group, T a maximal torus of G. For \(g \in G\) , denote by \(\tau(g)\) the automorphism of the manifold G/T induced by left multiplication by g.

\begin{enumerate}
\def\labelenumi{\alph{enumi})}
\item
  Let \(t\) be an element of \(\mathrm{T}\) such that the subgroup generated by \(t\) is dense in \(\mathrm{T}\) . Show that the fixed points of \(\tau(t)\) are the classes \(n\mathrm{T}\) for \(n \in \mathrm{N}_{\mathrm{G}}(\mathrm{T})\) . Deduce that \((\mathrm{N}_{\mathrm{G}}(\mathrm{T}):\mathrm{T}) = \sum_{i}(-1)^{i}\dim_{\mathbf{R}}\mathrm{H}^{i}(\mathrm{G / T},\mathbf{R})\) .
\item
  Let g be an arbitrary element of G; show that the set of fixed points of \(\tau(g)\) is non-empty, and deduce another proof of Th. 2.*
\end{enumerate}

\begin{enumerate}
\def\labelenumi{\arabic{enumi})}
\setcounter{enumi}{10}
\tightlist
\item
  Let G be a compact Lie group. A subgroup S of G is said to be of type (C) if it is equal to the closure of a cyclic subgroup that is of finite index in its normalizer.
\end{enumerate}

\begin{enumerate}
\def\labelenumi{\alph{enumi})}
\item
  Let \(S\) be a subgroup of type (C); show that \(S_0\) is a maximal torus of \(G_0\) , and that \(S\) is the direct product of \(S_0\) and a finite cyclic subgroup.
\item
  Prove that every element g of G is contained in a subgroup of type (C) (consider the group generated by g and a maximal torus of \(Z(g)_{0}\) ).
\item
  Let S be a subgroup of type (C), and s an element of S whose set of powers is dense in S. Show that every element of \(sG_{0}\) is conjugate under \(\operatorname{Int}(G_{0})\) to an element of \(sS_{0}\) (use the method of Exerc. 10).
\item
  Denote by \(p: G \to G/G_{0}\) the quotient map. Show that the map \(S \mapsto p(S)\) induces a bijection from the set of conjugacy classes of subgroups of G of type (C) to the set of conjugacy classes of cyclic subgroups of \(G/G_{0}\) .
\item
  Let S be a subgroup of G of type (C); denote by \(S_{\rho}\) the set of elements of S whose class in \(S/S_{0}\) generates \(S/S_{0}\) . Show that two elements of \(S_{\rho}\) that are conjugate in G are conjugate in \(N_{G}(S)\) .
\end{enumerate}